\documentclass{article}
	\usepackage[UTF8]{ctex}
	\usepackage{amsmath}
	\usepackage{graphicx}
	\usepackage{underscore}
	\usepackage{geometry}
	\usepackage{anysize}
	\date{}
	\title{三方坐镇题解}
	\geometry{top=10mm,bottom=15mm,left=25mm,right=25mm}
	
	\begin{document}
		\maketitle \vspace{-50pt}
		
		\section{注意}
			\indent 这道题非常容易写错！写的容易一定要细心，对代码进行检查时请多在纸上编写几个样例。\\
			\indent 另外，本题无法使用随机撒点方法，但是可适用于验证代码是否正确。
		
		\section{突破点}
			\indent 按照圆之间关系划分种类的方法是不可取的，因为数量太多了！\\
			\indent 由于题目难度降低，可以通过''若三个圆两两相交，则输出 -1 ''对本道题进行突破。
			
		\section{做法}
			\subsection{判断两个圆两两相交的情况}
				\indent 两个圆相交的条件是两圆距离大于两圆半径差值，小于两圆半径之和，而并不是大于等于，小于等于，这两种情况分别对应内切和外切。
			
			\subsection{去除被包含的圆的计算}
				\indent 圆A被圆B包含的条件是两圆距离小于等于圆A半径减去圆B半径，不能没有等于。\\
				\indent 计算没有被另外的圆包含的面积之和。\\
				\indent 若两(三)圆重合，保留一个圆。两圆重合的条件是圆心和半径相等。
				
			\subsection{去除两圆相交部分的计算}
				
				\indent 两圆相交部分只能计算一次。\\
				\indent 相交部分为圆1扇形+圆2扇形-两个三角形。三角形可以通过海伦定理求出，两个扇形的周角可以通过余弦定理求出。其它方法也可以，不过略为复杂。不推荐使用设三角形底边的高为参数，求解方程的方法，因为三角形为锐角三角形，直角三角形，钝角三角形的情况都存在。
\begin{figure}
\begin{minipage}[t]{0.5\linewidth}
\centering
\includegraphics[width=3.1in]{3.png}
\end{minipage}%
\begin{minipage}[t]{0.5\linewidth}
\centering
\includegraphics[width=3.1in]{4.png}
\end{minipage}
\end{figure}
				
	\section{本题延伸}
		\indent 关于三个圆的不重合面积之和(包括三个圆两两相交的情况)，可以采用圆的面积并(自行网络搜索)，这个方法也适用于多个圆的不重合面积之和。
	\end
	{document}
	